\section{System Design}
\label{sec:system-design}

\begin{figure}[t]
  \centering
  \begin{tikzpicture}[
  node distance=4.5mm,
  box/.style={
    draw=black,
    rounded corners=1.5pt,
    align=center,
    inner sep=5pt,
    text width=0.78\linewidth,
    font=\small
  },
  arrow/.style={-{Stealth[length=2.2mm]},thick}
]
  \node[box] (request) {HTTP request: text, VoiceDesign description, language, seed, and output mode};
  \node[box,below=of request] (server) {Rust HTTP server: validation, request identity, admission, cancellation, and response framing};
  \node[box,below=of server] (scheduler) {Fixed-capacity Rust scheduler: request-local state, bounded PCM rings, and backpressure};
  \node[box,below=of scheduler] (talker) {VoiceDesign talker and 15-step code predictor on CUDA/cuBLAS};
  \node[box,below=of talker] (handoff) {CUDA-event-ordered device-to-device handoff of 16 codec codes per frame};
  \node[box,below=of handoff] (decoder) {Stateful native neural speech decoder on CUDA/cuBLAS};
  \node[box,below=of decoder] (response) {Progressive 24 kHz PCM multipart stream or bounded buffered WAV response};

  \draw[arrow] (request) -- (server);
  \draw[arrow] (server) -- (scheduler);
  \draw[arrow] (scheduler) -- (talker);
  \draw[arrow] (talker) -- (handoff);
  \draw[arrow] (handoff) -- (decoder);
  \draw[arrow] (decoder) -- (response);
\end{tikzpicture}

  \caption{Implemented online path. Immutable model weights are shared, while
  mutable request state and CUDA execution resources are request-local. The
  diagram is a vector TikZ source and encodes no benchmark result.}
  \label{fig:architecture}
\end{figure}

Figure~\ref{fig:architecture} shows the complete online path. A single process
owns HTTP validation, scheduling, model execution, neural decoding, and
response delivery. The server loads one shared engine, creates a fixed-capacity
scheduler, and completes a real one-frame native warm-up before binding its
listener. Readiness therefore establishes that the model and end-to-end native
pipeline have produced a valid PCM packet; it is not merely a process-liveness
signal.

\subsection{Admission and ownership}

The canonical synthesis endpoint validates bounded JSON before attempting
engine admission. Text, textual VoiceDesign instruction, language, duration,
sampling settings, and response mode are checked at this boundary. Each
admitted request receives a unique identity and owns its cancellation state,
generation session, decoder session, CUDA stream and events, cuBLAS handle,
random state, PCM buffers, and scheduler slot. Immutable model weights are
shared across sessions.

The scheduler supports one through six active slots. Pending starts can be
coalesced for at most 2 ms, but capacity is fixed and exhaustion is explicit.
Each request has three reusable PCM buffers. A request advances only when a
buffer and queue position are available, so a slow client applies backpressure
instead of causing unbounded audio allocation.

\subsection{Packet cadence}

One codec frame represents 1,920 output samples, or 80 ms at 24 kHz. The first
decoder step requests one frame to minimize first-audio latency. Subsequent
steps request up to four frames and therefore produce at most 7,680 samples per
packet. This policy bounds buffering without forcing every later step to pay
the overhead of a one-frame packet.

\subsection{Lifecycle and failure handling}

Natural codec end-of-sequence closes a request after all owned PCM has been
delivered. Cancellation is propagated through the scheduler to native state,
and resource retirement is bounded. Model-contract, ABI, CUDA, packet-order,
and state-transition failures are surfaced instead of silently falling back to
a different implementation. The process does not attempt unsafe in-process
recovery of a stuck CUDA session.
